% 张三慧《大学物理学》第4版第一卷第6章《振动》
% 原教材PDF物理页209--212（印刷193--196页）。
% 章末习题6.1--6.27；不可与思考题6.1--6.10混淆。
% 图6.21、6.22--6.27均需要后续矢量重绘与独立图示核验。
\section*{第6章\quad 振动}
\subsection*{章末习题6.1--6.27}

\paragraph{习题 6.1}
一个小球和轻弹簧组成的系统，按
\[
 x=0.05\cos(8\pi t+\pi/3)
\]
的规律振动。
（1）求角频率、周期、振幅、初相、
最大速度及最大加速度；
（2）求 \(t=1\,\mathrm s,2\,\mathrm s,10\,\mathrm s\)
等时刻的相；
（3）分别画出位移、速度、加速度与时间的关系曲线。

\paragraph{习题 6.2}
有一个和轻弹簧相连的小球，沿 \(x\) 轴作振幅为
\(A\) 的简谐运动，振动的表达式用余弦函数表示。
若 \(t=0\) 时的运动状态分别为：
（1）\(x_0=-A\)；
（2）过平衡位置向 \(x\) 正方向运动；
（3）过 \(x=-A/2\) 处，且向 \(x\) 负方向运动。
试用相量图法分别确定相应的初相。

\paragraph{习题 6.3}
已知一个谐振子（即作简谐运动的质点）的
振动曲线如原书图6.21所示。
（1）求 \(a,b,c,d,e\) 各状态相应的相；
（2）写出振动表达式；
（3）画出相量图。
\textit{原图可读信息：}纵轴 \(x/\mathrm{cm}\)，
振幅 \(5.0\,\mathrm{cm}\)，\(t=0\) 时
\(x=+2.5\,\mathrm{cm}\) 且上升；
图上标出 \(t=1.0\,\mathrm s\) 时
零位向下穿过，以及 \(t=2.2\,\mathrm s\)
时回到 零点且上升。
\(a\) 是正极大，\(b\) 是 \(x=2.5\) 的下降支，
\(c\) 是下行零点，\(d\) 是 \(x=-2.5\) 的下降支，
\(e\) 是 \(x=-2.5\) 的上升支。

\paragraph{习题 6.4}
作简谐运动的小球，速度最大值为
\(v_m=3\,\mathrm{cm/s}\)，振幅
\(A=2\,\mathrm{cm}\)。
若从速度为正的最大值的某时刻开始计算时间，
（1）求振动周期；
（2）求加速度的最大值；
（3）写出振动表达式。

\paragraph{习题 6.5}
一水平弹簧振子，振幅
\(A=2.0\times10^{-2}\,\mathrm m\)，
周期 \(T=0.50\,\mathrm s\)。
当 \(t=0\) 时：
（1）振子过 \(x=1.0\times10^{-2}\,\mathrm m\)
处，向负方向运动；
（2）振子过 \(x=-1.0\times10^{-2}\,\mathrm m\)
处，向正方向运动。
分别写出两种情况下的振动表达式。

\paragraph{习题 6.6}
两个谐振子作同频率、同振幅的简谐运动。
第一个的振动表达式为
\(x_1=A\cos(\omega t+\varphi)\)。
当第一个振子从振动正方向回到平衡位置时，
第二个振子恰在正方向位移的端点。
（1）求第二个振子的振动表达式和二者的相差；
（2）若 \(t=0\) 时
\(x_1=-A/2\)，且向 \(x\) 方向运动，
画出二者的 \(x-t\) 曲线及相量图。
\textit{扫描歧义：}原书第（2）问的
「向 \(x\) 方向运动」未明确写正或负号，
需对其作条件性说明。

\paragraph{习题 6.7}
两个质点平行于同一直线并作同频率、同振幅的
简谐运动。
在振动过程中，每当它们经过振幅一半的地方时相遇，
而运动方向相反。
求它们的相差，并作相量图表示。

\paragraph{习题 6.8}
一弹簧振子，弹簧劲度系数
\(k=25\,\mathrm{N/m}\)。
当振子以初动能 \(0.2\,\mathrm J\)
和初势能 \(0.6\,\mathrm J\) 振动时：
（1）振幅多大？
（2）位移多大时势能和动能相等？
（3）位移多大时动能与势能之比最大？

\paragraph{习题 6.9}
将一劲度系数为 \(k\) 的轻质弹簧
上端固定悬挂，下端挂一质量为 \(m\) 的小球，
平衡时弹簧伸长为 \(b\)。
试写出以此平衡位置为原点的小球动力学方程，
证明小球将作简谐运动并求其周期。
若振幅为 \(A\)，总能量是否是
\(\tfrac12kA^2\)？
（总能量包括动能、重力势能以及弹簧的弹性势能，
两种势能均取平衡位置为零。）

\paragraph{习题 6.10}
在分析图6.6所示弹簧振子时，
通常忽略弹簧质量。
现考虑弹簧质量的影响。
设弹簧质量为 \(m'\)，沿弹簧长度均匀分布，
振子质量为 \(m\)。
用 \(v\) 表示振子某时刻的速度，
弹簧各点速度与到固定端的长度成正比。
证明：
（1）此刻弹簧振子的动能为
\[
 \frac12\left(m+\frac{m'}3\right)v^2,
\]
因此系统有效质量为 \(m+m'/3\)；
（2）系统角频率应为
\[
 \omega=\sqrt{\frac{k}{m+m'/3}}.
\]

\paragraph{习题 6.11}
将劲度系数分别为 \(k_1,k_2\) 的
两根轻弹簧串联在一起，竖直悬挂，
下面系一质量为 \(m\) 的物体，
做成一个在竖直方向振动的弹簧振子。
试求其振动周期。

\paragraph{习题 6.12}
劲度系数分别为 \(k_1,k_2\) 的
两根轻弹簧和质量为 \(m\) 的物体相连，
如图6.22所示，试写出物体的动力学方程
并证明该振动系统的周期为
\[
 T=2\pi\sqrt{\frac m{k_1+k_2}}.
\]
\textit{图示：}滑块左右分别连接固定墙上的两弹簧，
滑块水平无摩擦振动。

\paragraph{习题 6.13}
在水平光滑桌面上用轻弹簧连接
两个质量都是 \(0.05\,\mathrm{kg}\) 的小球
（图6.23），弹簧劲度系数
\(1\times10^3\,\mathrm{N/m}\)。
今沿弹簧轴线向相反方向拉开两球然后释放，
求此后两球振动的频率。

\paragraph{习题 6.14}
设想穿过地球挖一条直细隧道（图6.24），
隧道壁光滑。
在隧道内放一质量为 \(m\) 的小球，
它离隧道中点的距离为 \(x\)。
设地球为均匀球体，质量为 \(M_E\)，
半径为 \(R_E\)。
（1）求球受到的重力
（提示：小球只受球所在处的圆球面以内的地球质量的引力）。
（2）证明球在隧道内在重力作用下作简谐运动，
并求周期。
（3）近地圆轨道人造地球卫星的周期多大？

\paragraph{习题 6.15}
一物体放在水平木板上，
物体与板间的静摩擦系数为 \(0.50\)。
（1）木板沿水平方向作频率为 \(2.0\,\mathrm{Hz}\)
的简谐运动时，要使物体在板上不致滑动，
振幅最大应是多少？
（2）木板沿竖直方向作简谐运动，
振幅为 \(5.0\,\mathrm{cm}\)，
要使物体一直保持与板面接触，
振动频率最大应是多少？

\paragraph{习题 6.16}
如图6.25，一块均匀长木板质量为 \(m\)，
对称地平放在相距 \(l=20\,\mathrm{cm}\)
的两个滚轴上。
两滚轴的转动方向相反，
已知滚轴表面与木板间的摩擦系数为
\(\mu_k=0.5\)。
今使木板沿水平方向移动一段距离后释放，
证明木板将作简谐运动并求其周期。
\textit{图示：}两个滚轴位于木板质心左右
\(l/2\) 处，旋转方向使两接触点的摩擦力
分别指向系统中心。

\paragraph{习题 6.17}
质量为 \(m=121\,\mathrm g\) 的水银
装在U形管中，管截面积
\(S=0.30\,\mathrm{cm^2}\)。
当水银面上下振动时，其振动周期
\(T\) 是多大？
水银密度为
\(13.6\,\mathrm{g/cm^3}\)。
忽略水银与管壁的摩擦。

\paragraph{习题 6.18}
行星绕太阳的运行速度可分解为径向分量
\(v_r=dr/dt\) 和角向分量
\(v_\theta=r\,d\theta/dt\)。
因此质量 \(m\) 的行星机械能为
\[
 E=\frac12m\left(\frac{dr}{dt}\right)^2
 +\frac12mr^2\left(\frac{d\theta}{dt}\right)^2
 -\frac{GmM_S}{r}.
\]
（1）证明：若 \(L\) 是行星对太阳的
恒定角动量，则
\[
 E=\frac12m\dot r^2+\frac{L^2}{2mr^2}
 -\frac{GmM_S}{r}.
\]
（2）圆轨道 \(r=r_0\) 与近似圆轨道
\(r=r_0+x\)，\(x\ll r_0\)；
证明近圆轨道有
\[
 E=\frac12m\dot x^2+
 \frac{3L^2}{2mr_0^4}x^2
 -\frac{GmM_S}{r_0^3}x^2
 -\frac{GmM_S}{2r_0}.
\]
（3）与简谐振动能量比较，
证明径向振动的等效劲度系数
\(k=GM_Sm/r_0^3\)。
（4）证明径向振动周期等于行星绕圆轨道公转的周期，
并画近圆轨道的轨迹。

\paragraph{习题 6.19}
一质量为 \(m\) 的刚体在重力矩作用下，
绕固定的水平轴 \(O\) 作小幅度无阻尼摆动，
如图6.26。
刚体质心 \(C\) 到轴 \(O\) 的距离为 \(b\)，
刚体对轴 \(O\) 的转动惯量为 \(I\)。
试用转动定律写出动力学方程，
证明 \(OC\) 与竖直线的夹角
\(\theta\) 的变化为简谐运动，
且周期
\[
 T=2\pi\sqrt{\frac I{mgb}}.
\]

\paragraph{习题 6.20}
一细圆环质量为 \(m\)，半径为 \(R\)，
挂在墙上的钉子上。
求它的微小摆动周期。

\paragraph{习题 6.21}
HCl 分子中两离子的平衡间距为
\(1.3\times10^{-10}\,\mathrm m\)，
势能可近似表示为
\[
 E_p(r)=-\frac{e^2}{4\pi\varepsilon_0r}+\frac{B}{r^9},
\]
其中 \(r\) 为两离子间距离。
（1）求 HCl 分子的微小振动频率。
由于Cl离子的质量比质子质量大很多，
可以认为Cl离子不动。
（2）利用式（6.30），并设
HCl分子处于基态振动能级 \(n=1\)，
按经典简谐运动计算，求其中质子振动的振幅。

\paragraph{习题 6.22}
一单摆在空气中摆动，摆长
\(1.00\,\mathrm m\)，初始振幅角
\(\theta_0=5^\circ\)。
经过 \(100\,\mathrm s\)，振幅减为
\(\theta_1=4^\circ\)。
再过多久，其振幅减为
\(\theta_2=2^\circ\)？
此单摆的阻尼系数多大？品质因数 \(Q\) 多大？

\paragraph{习题 6.23}
证明：当驱动力的频率等于系统的
固有频率时，受迫振动的速度幅达到最大值。

\paragraph{习题 6.24}
一质点同时参与两个在同一直线上的简谐运动，
其表达式为
\[
 x_1=0.04\cos(2t+\pi/6),\qquad
 x_2=0.03\cos(2t-\pi/6).
\]
试写出合振动的表达式。

\paragraph{习题 6.25}
三个同方向、同频率的简谐振动为
\[
 x_1=0.08\cos(314t+\pi/6),\quad
 x_2=0.08\cos(314t+\pi/2),\quad
 x_3=0.08\cos(314t+5\pi/6).
\]
求：（1）合振动的角频率、振幅、初相和振动表达式；
（2）合振动由初始位置运动到
\(x=(\sqrt2/2)A\)（\(A\) 为合振幅）
所需最短时间。

\paragraph{习题 6.26}
一质点同时参与相互垂直的两个简谐运动：
\[
 x=0.06\cos(20\pi t),\qquad
 y=0.04\cos(20\pi t+\pi/2).
\]
证明其轨迹为正椭圆（长短轴分别与两坐标轴重合），
求长半轴、短半轴及绕行周期。
质点绕行是右旋（顺时针）还是左旋（逆时针）？

\paragraph{习题 6.27}
李萨如图可用来测量频率。
示波器的水平和垂直输入端分别加上交变电压，
其荧光屏出现原书图6.27所示闭合曲线。
已知水平方向振动的频率为
\(2.70\times10^4\,\mathrm{Hz}\)，
求垂直方向的振动频率。
\textit{图示核验：}图6.27呈复杂的有交叉点的闭合
李萨如曲线，横纵瓣数需依原图独立计数，不能仅凭文字猜比值。
